\part{2026}
\section{
    0407联合早报·美国阿耳忒弥斯2号打破人类距地球最远飞行纪录
} {
    \setlength{\parindent}{2em}
    （休斯敦综合电）美国东部时间星期一（4月6日）18时40分许（新加坡时间7日6时40分许），执行美国阿耳忒弥斯2号载人绕月飞行任务的猎户座飞船处于月球背面，四名宇航员进入预定的约40分钟通信中断状态。

美国国家航空航天局（NASA）说，绕月飞行的阿耳忒弥斯2号（Artemis 2）星期一13时57分，打破了1970年阿波罗13号月球任务创造的人类距离地球最远飞行纪录。

阿波罗13号曾飞至距离地球约40万171公里处。这次阿耳忒弥斯2号最远飞至距离地球约40万6771公里，较原纪录超出6600公里。

专家表示，飞得远并非阿耳忒弥斯2号最终目的，验证载人地月往返技术与安全体系才是其核心任务。这是实现人类长期重返月球、开展未来火星任务的关键一步。

美国时隔半个多世纪再次送人赴月球，此次绕月飞行历时六小时，搭载四名宇航员的猎户座飞船掠过月球表面，最近时距离月面约6400公里。

宇航员全程紧盯着窗外，仔细描述了月面特征。他们看到日食，即月球运行到太阳前方，也看到流星体撞击月面，迸发出闪光。

在飞往月球背面途中，宇航员为一些此前没有命名的月球地貌赋予了新的临时名称。在发给休斯顿任务控制中心的无线电信息中，宇航员汉森建议将其中一个陨石坑命名为“诚信”（Integrity），与猎户座飞船同名；另一个陨石坑则以宇航员怀斯曼亡妻卡罗尔（Carroll）的名字命名。

汉森后来说，宇航员观测到一些月球地貌，“人类从没见过，即使是阿波罗计划也没见过”。

由于月球自转和绕地球公转的速度相同，它的背面始终背对着地球，目前只有阿波罗和阿耳忒弥斯的宇航员曾直接观测过月球背面。

阿耳忒弥斯2号任务首席科学家杨格对宇航员说：“今天你们让我们感觉离月球更近了，我们感激不尽。”

绕月飞行完成后，美国总统特朗普与宇航员连线，简短通话。特朗普祝贺宇航员创造了历史，并说他们“让整个美国感到非常自豪，无比自豪” 。

特朗普一直向NASA施压，要求在他的第二个任期2029年1月结束前，再次实现载人登月。

source: https://www.zaobao.com/news/world/story20260407-8856596
}